\chapter[Der am Kreuz]{Der am Kreuz ist meine Liebe}
\quellseite{8}

Dieses schöne Passionslied (No. 71) ist von Joh. Ernst Greding († 1748 als Pfarrer zu Allheim bei Hanau) gedichtet worden nach dem Wahlspruch des frommen Märtyrers Ignatius: „Meines Herzens Liebe ist am Kreuze gestorben.“ Diesen Wahlspruch hat Ignatius, der wegen seines christlichen Bekenntnisses hochangesehene Bischof von Antiochien in Syrien, bis zu seinem Tode festgehalten. Als er nämlich um das Jahr 107 in Rom den wilden Thieren vorgeworfen wurde, hat er auf dem Wege zur Marter immer die Worte gesagt: „Ei, du süßer Name Jesus!“ Und auf die Frage, weshalb er immer diese Worte wiederhole, hat er geantwortet: „Dieser Name ist mir ins Herz hineingegraben, ich kann ihn nicht vergessen, und so oft ich ihn nenne, fühlt mein Herz neuen Trost.“

Der bekannte fromme Liederdichter Valerius Herberger hat diese Geschichte vom heiligen Ignatius sehr gern gehört. In gleicher Liebe zu dem gekreuzigten Gottessohn singt er in seinem schönen Sterbelied (No. 426):

\begin{verse}
In meines Herzens Grunde\\
Dein Nam und Kreuz allein\\
Funkelt all Zeit und Stunde,\\
Drauf kann ich fröhlich sein.\\
Erschein mir in dem Bilde\\
Zum Trost in meiner Noth,\\
Wie du, Herr Christ, so milde\\
Dich hast geblut zu Tod!
\end{verse}

Auch Greding, der Verfasser unsers Liedes, bezeuget mit seinem Gesang, daß ihm Jesus der Gekreuzigte über alles ging. Obwohl er mit dem kaiserlichen Dichterlorbeer geschmückt war, so hat er doch um „den ewigen Kranz“ sein armes Leben dem Herrn Christo ganz ergeben. Und wie von dem Leibe des gläubigen Zeugen Christi Ströme lebendigen Wassers fließen, Joh. 7, 38., so hat auch dieses Lied in manches Menschen Herzen die Liebe Christi erweckt oder gestärkt.

Am Ende des 18. Jahrhunderts wurde der fromme Minister v. Seckendorf von Herzog Karl von Württemberg zu einem Hofball befohlen. Der Minister erbat sich die Gnade, das erste Musikstück auswählen zu dürfen, und da die Bitte gewährt wurde, so ertönte, als der Tanz beginnen sollte, durch den Saal die Weise unsers Liedes:

\begin{verse}
Der am Kreuz ist meine Liebe,\\
Meine Lieb ist Jesus Christ.\\
Weg, ihr argen Seelendiebe,\\
Satan, Welt und Fleischeslist!\\
Eure Lieb ist nicht von Gott,\\
Eure Lieb ist gar der Tod;\\
Der am Kreuz ist meine Liebe,\\
Weil ich mich im Glauben übe.

Der am Kreuz ist meine Liebe:\\
Freuler, was befremdet dich,\\
Daß ich mich im Glauben übe?\\
Jesus gab sich selbst für mich:\\
So ward er mein Friedensschild,\\
Aber auch mein Lebensbild.\\
Der am Kreuz ist meine Liebe,\\
Weil ich mich im Glauben übe.
\end{verse}

Der Herzog forderte seinen Minister nicht mehr zum Tanze auf, und dieser hatte hinfort Ruhe vor dergleichen Zumuthungen, die Liebe zu Christo zu verlezen.

Abgeschreckt von der Weltliebe und erweckt zur Liebe gegen den Gekreuzigten wurde im Jahre 1775 ein 17jähriger Jüngling in Altdorf, als man in der Charwoche dieses Lied in der Kirche sang, besonders durch den dritten und vierten Vers:

\begin{verse}
Der am Kreuz ist meine Liebe:\\
Sünde, du verlierst den Sturm.\\
Weh mir, wenn ich den betrübe,\\
Der statt meiner ward ein Wurm!\\
Kreuzigt ich nicht Gottes Sohn?\\
Trät ich nicht sein Blut mit Hohn?\\
Der am Kreuz ist meine Liebe,\\
Weil ich mich im Glauben übe.

Der am Kreuz ist meine Liebe:\\
Schweig, Gewissen! niemand mahnt;\\
Gott preist seine Liebestrebe,\\
Wenn mir von der Handschrift ahnt.\\
Schau, wie mein Halsbürge zahlt;\\
Gottes Blut hat sie durchmalt.\\
Der am Kreuz ist meine Liebe,\\
Weil ich mich im Glauben übe.
\end{verse}

\quellseite{9}

Diese Worte drangen dem Jüngling so ins Herz, daß er von Stund an aller Lustbarkeit der Welt entsagte und lieber sterben wollte, als den noch einmal betrüben, der für uns den bittern Tod des Kreuzes aus Liebe erduldet hat. Er hat dann später selbst Jesu-Lieder gedichtet.

Auch der fünfte Vers hat mindestens Einen Menschen vor Verarmung bewahrt und in der Treue gegen den Gekreuzigten gestärkt. Er lautet so:

\begin{verse}
Der am Kreuz ist meine Liebe:\\
Drum, Tyrann, nun foltre, stoß!\\
Hunger, Blöße, Henkerhiebe,\\
Nichts macht mich von Jesu los,\\
Nicht Gewalt, nicht Gold, nicht Ruhm,\\
Engel nicht, kein Fürstenthum.\\
Der am Kreuz ist meine Liebe,\\
Weil ich mich im Glauben übe.
\end{verse}

An einem Abend des Jahres 1815 nämlich wurde der deutsche Jüngling Joh. Kullen in Frankreich durch die schlüpfrigen Gespräche seiner Stubengenossen am Einschlafen gehindert. Um womöglich ihre Worte gar nicht zu hören, sagte er früher auswendig gelernte Bibelsprüche und Gesänge für sich her, und so kam ihm auch unser Lied in den Sinn. Dabei wurde es ihm klar, daß er in der Liebe zu seinem Heiland kälter geworden war. Eine Stimme in ihm sagte: „Die Worte des Liedes gehen dich an.“ Er dagegen fragte sich: „Was willst du denn? Willst du ein sündliches Leben führen wie diese? Möchtest du Reichthum haben in der Welt oder Ehre? Möchtest du ein Officier sein?“ Und er kam durch Gottes Gnade schließlich zu der Antwort: „Ich möchte nur dich haben, lieber Heiland!“ „Nun“, hieß es wieder in ihm, „so geht es dich an: ‚Der am Kreuz ist meine Liebe.‘“ Und der Gekreuzigte mit seiner Liebe behielt den Sieg im Herzen dieses Mannes. Er entsagte weltlichen Ehren, Gütern und Genüssen.

Der letzte Vers endlich sollte sich besonders an dem Dichter unseres Liedes selber bewähren. Nachdem dieser noch am Gründonnerstag und Charfreitag 1748 das Wort von Jesu dem Gekreuzigten mit heiligem Liebesfeuer verkündigt hatte, ging er schnell und leicht, Tags darnach, am Ostermontag, durch einen seligen Tod „ins himmlische Vaterhaus“.

Möge Jesus der Gekreuzigte immer mehr auch unsers Herzens Liebe werden, auf daß er auch einst unser Gewinn sei im Sterben und so auch an uns dann sich bewähre der Schlußvers unsers Liedes:

\begin{verse}
Der am Kreuz ist meine Liebe:\\
Komm, Tod, komm, mein bester Freund!\\
Wenn ich wie ein Staub zerstiebe,\\
Wird mein Jesus mir vereint;\\
Da, da schau ich Gottes Lamm,\\
Meiner Seelen Bräutigam.\\
Der am Kreuz ist meine Liebe,\\
Weil ich mich im Glauben übe.
\end{verse}
